(1) 自相似性给 \(Z=R_s+[R_pZ/(R_p+Z)]\)，即 \(Z^2-R_sZ-R_sR_p=0\)。物理输入阻抗必须正，故 \(Z=[R_s+\sqrt{R_s^2+4R_sR_p}]/2\)，负根不可作为被动电路阻抗。

(2) 输入电流 \(I_0=V/Z\)，串联 \(R_s\) 后节点电势 \(V_1=V-I_0R_s=V(Z-R_s)/Z\)。首并联支路 \(I_p=V_1/R_p\)，后续无限子网络 \(I_{\rm tail}=V_1/Z\)，且因 \(1/(Z-R_s)=1/R_p+1/Z\)，有 \(I_p+I_{\rm tail}=I_0\)。

(3) 把最后一级后续断开取 \(Z_1=R_s+R_p\)，迭代 \(Z_{n+1}=R_s+R_pZ_n/(R_p+Z_n)\)。映射 \(f'(z)=R_p^2/(R_p+z)^2\in(0,1)\)，故它是严格收缩的正值映射，\(Z_n\) 由高于固定点的初值单调下降至 \(Z\)。也可将边界改为其它非负负载，极限唯一。

(4) 首级电功率 \(P_s=I_0^2R_s\)，并联功率 \(P_p=V_1^2/R_p\)，后续网络功率 \(P_{\rm tail}=V_1^2/Z\)。三者之和 \(I_0^2R_s+V_1^2(1/R_p+1/Z)=I_0^2[R_s+(Z-R_s)]=I_0^2Z=V^2/Z\)。当 \(R_p\to\infty\)，\(Z\sim\sqrt{R_sR_p}\to\infty\)，无泄漏回路等效断路；\(R_p\to0\) 时 \(Z\to R_s\)，首节点直接短路。
